% ===========================================================================
% OWNER: M1 (Zoya Ahmad)
% GRADED BY: Form 2 #6 / #8 individual ownership
% DRAW FROM: git log --author, your merged PRs, docs/09-Team/members.md, evidence/
% Rules: mid-eval-report/CLAUDE.md. Do not edit a file you do not own.
% ===========================================================================
\subsection{Member 1: Zoya Ahmad}
\textbf{What I built.} I own the \path{evidence/} package: the static resource-evidence extractor
that names the shared resource a polluter writes and a victim reads, with the exact bytecode
locations (W8). I wrote its contract (\path{docs/contracts/resource-evidence.md}, \pr{8}), then built
it in phases: lifecycle attribution at depth~1 (\pr{9}), configurable call depth (\pr{12}),
polluter-to-victim edges (\pr{15}), an automatic check against the fixture's ground truth
(\pr{23}), the first real-project run on fastjson (\pr{24}), lifecycle writes as polluter edges
(\pr{28}) and auto-deepening (\pr{34}). I proposed ADR-002 and ADR-006 and wrote the panel-action
note comparing one-by-one search with delta debugging (A6, \pr{50}). For this report I wrote
Chapters~1 and~2, Sections~5.6--5.8 with their three diagrams, and my module section
(Section~\ref{ch6:sec:m1}). On \texttt{main} on 12 October 2026 that is 23 merged pull requests and
56 commits under my name.

\textbf{Reviews.} I reviewed 36 pull requests written by Member~2 or Member~3. My review comments on
GitHub are short approvals; the checks behind them are in \path{docs/evidence-m1.md}. For example,
before approving Member~2's \pr{49} (the runner side of ADR-008), I ran all three test suites on its
commit \texttt{fa621af} (evidence 39, evaluation 80 and runner 144 tests, all passing), confirmed
that it did not change how \path{runner/cli.py} calls my \texttt{analyse\_pair} and
\texttt{report\_fields}, and ran the real \texttt{diagnose} command on fixture case F4, which ended
\texttt{VERIFIED} with my module's edge on \texttt{LateFlag.isSet}.

\textbf{One decision: deepen only when needed (ADR-006).} ADR-002 fixed the analysis depth at 2,
with at most 3. When I measured the two real fastjson cases, depths 1--3 found none of the 42 tests
that really affect FJ-01's victim; depth~4 found all 42, and FJ-02 needed depth~5
(Table~\ref{ch6:tab:m1-depth}). I considered keeping 1--3 (honest, but the tool explains neither
real case), raising the default to~5 (every pair pays the deepest cost, 26.3~s for one FJ-02 pair,
and the most over-approximation), and allowing 1--5 with deeper analysis only on request (but the
\texttt{diagnose} command always uses the default, so it would still find nothing). I chose
auto-deepening with the range 1--5: start at~2 and go one level deeper only when no edge was found
and the walk was cut off. Because this changed the contract, Members~2 and~3 had to agree, and they
did (\pr{29}, \pr{33}). The cost is a repeated analysis for each extra level, and a limitation line
in the report whenever the evidence is deeper than the default.

\textbf{One limitation of my module.} The extractor sees \emph{that} a test writes a field, not
\emph{what} it writes. A test that changes a value and restores it in \texttt{tearDown} still
counts as a writer: on FJ-01 this produced one false edge, \texttt{DateTest2\#test\_date}. For the
same reason it cannot tell which of two written fields causes a failure, as with fastjson's time
zone and locale. FlakeTrace never blames a test on such an edge alone, because the polluter must
also be confirmed by reruns, but the explanation can still name the wrong field. Fixing it needs
runtime evidence on the already-reduced order, the fallback planned in ADR-001; I have not built it
yet.
